% Chapter 5

\chapter{总结与展望}

\section{工作总结}

本文围绕基于LoRA技术的特定风格视频生成模型展开研究。针对当前视频生成领域普遍存在的训练成本高、特定风格定制困难以及时序一致性难以保持等问题，本文提出并验证了一条从图像风格学习到短视频原型验证的技术路线。

本文系统分析了扩散模型、潜空间扩散框架以及LoRA的基本原理，明确了在有限算力资源下，利用LoRA对U-Net注意力模块进行局部重塑，是实现特定风格迁移的一条高效途径。同时，结合AnimateDiff等时序建模思想，本文确立了先进行外观学习、再开展时序验证的递进式研究策略。

针对目标动画风格，本文构建了包含810张图像和5组视频片段的实验数据集。预处理阶段主要完成了中心裁剪、尺寸统一和人工校验等工作，以尽量保证训练样本在风格上的一致性；同时为图像样本配套整理了成对的caption，为模型提供相对稳定的条件引导信号。

本文以Stable Diffusion 1.5为基础模型，通过比较300步短训与2000步长训的结果，验证了整条训练链路能够稳定运行。随后又引入StorybookRedmond和Counterfeit-V3.0等具有不同视觉先验的底模进行替换实验。实验结果说明，动画先验更强的底模，往往能在像素平均绝对差等指标上表现出更明显的风格迁移能力，但与此同时，也更容易把自身固有的审美倾向带入结果，并对原有构图关系造成改写。

由于显式视频训练的成本较高，本文采用递推式短视频生成方案进行了原型验证。量化结果表明，加载LoRA后的视频在MAD和MED指标上分别改善了约6.43\%和2.65\%，这说明图像阶段学到的风格能力确实能够在一定程度上延续到连续帧中。不过，LS指标仍增加了约11.72\%，这也说明当前方案在亮度稳定性方面还有比较明显的不足。

\section{研究展望}

尽管本文在低成本特定风格视频生成方面取得了一定进展，但受限于算力条件和实验周期，现有工作仍有不少可以继续深入的地方。后续研究可以从以下几个方向展开。

1. 引入显式的时序建模模块。当前视频原型主要依赖递推式生成，对运动轨迹和跨帧长程依赖的刻画还不够。后续可以考虑把训练好的风格LoRA与AnimateDiff运动模块，或更复杂的时空注意力机制进行更深入的结合，从生成机制层面缓解亮度闪烁和主体一致性不足的问题。

2. 增强结构可控性与泛化能力。实验中可以看到，风格注入较强时，有时会破坏原始构图。后续可引入ControlNet等条件控制技术，利用边缘图或深度图对生成过程施加几何约束，在保持风格特征的同时，更好地保留主体结构和镜头关系。

3. 完善多维度评估体系。目前使用的MAD、MED、LS等指标主要反映像素层面的稳定性。后续还应引入更贴近视觉感知的评价标准，例如衡量图文匹配程度的CLIP Score、衡量视频感知质量的FVD，以及更细粒度的风格相似度人工评测方案，从而形成更完整的视频生成评价体系。
